\section{From policy documents to ODRL: what must survive translation}
\label{sec:odrl-review}

The next proposed output of LegalMath is a machine-readable policy in the
Open Digital Rights Language, ODRL. The immediate case is the joint
SFC--HKMA circular of 28 July 2023, reference 23EC35, on the Streamlined
Approach for sophisticated professional investors (SPIs). The useful question
is whether an intermediary may apply particular streamlined procedures to a
client and transaction, and which duties remain. It is not simply whether a
trade is allowed. This difference determines what the translation must express
\citep{sfcspi,sfcspiannex1,sfcspiannex2}.

Consider a designated account whose gross exposure has risen above the client's
streamlining threshold. A sentence-level translator might read the threshold
rule and produce a prohibition on every further transaction. Annex~2 FAQ~6
expressly permits the account to continue operating and executing transactions
while top-ups and deposits are restricted. It does not expect the intermediary
to reduce or unwind exposure for this purpose. An ordinary suitability route
also remains relevant. The missing distinction is the kind of action and the
monitoring arrangement, not a different spelling of the threshold
\citep[FAQ~6]{sfcspiannex2}.

This review therefore recommends \emph{ODRL as the first new deliverable},
produced from a reviewed interpretation linked to the source. Selected Catala
calculations should accompany it. Requiring a complete Catala program first
would postpone the desired policy representation without resolving its
normative questions. Treating ODRL as the only representation of the circular
would instead risk losing meanings that the selected ODRL profile cannot
execute. The implementation proposal below develops a shared interpretation
with explicit projections into these different languages.

\subsection{Scope and strength of the evidence}

The review was completed on 22 October 2026. It combines the circular body,
both annexes, direct text-to-ODRL research, work on regulatory profiles and
formal semantics, and static inspection of selected software. The official
annex PDFs match the repository's archived copies byte for byte. This fixes a
historical translation case; a current operational application would still
need the relevant source editions, amendments and incorporated provisions.
The Code of Conduct, Professional Investor Rules and referenced FAQs must
remain named dependencies rather than silently assumed complete
\citep{hkmaodrlsurvey2026}.

Searches used Crossref, OpenAlex, GitHub, the W3C ODRL community's implementation
landscape, and references from the inspected papers. The recorded
``HKMA ODRL'' query returned no results, and no direct HKMA translation study
was found in the inspected material. This is a targeted survey, not proof
that no such work exists. Eight newly retained papers received technical
reading of the relevant method, semantics and evaluation; existing Catala
and eFLINT papers supplied focused comparisons. Nine public code checkouts
were pinned. Unavailable originals and documentation-only findings remain
identified in the accompanying reading record. No generation experiment or
third-party software benchmark was run \citep{hkmaodrlsurvey2026}.

The important evaluation question is whether a translation preserves the
decision and duty distinctions in the circular. A graph can be valid RDF,
conform to a set of shapes, and still prohibit the transactions that FAQ~6
allows. Conversely, a disagreement between evaluators can reflect different
declared semantics rather than a malformed graph. The following discussion
keeps these questions separate so that published scores are not asked to
answer a question their experiments did not test.

\subsection{ODRL's basic model and the need for a profile}

ODRL is a W3C model for expressing permissions, prohibitions and duties
concerning identifiable resources. An action says what may, may not, or must
be done; a party identifies a participant; an asset identifies the resource
to which the action applies. Constraints restrict when a rule applies.
The same policy can be serialized as JSON-LD, a JSON representation that
maps names to identifiers in an RDF graph. RDF supplies the linked statements;
it does not infer the legal meaning of ``sophisticated investor''
\citep[secs.~2--3]{odrlmodel2018}.

The distinction between two ODRL relationships is particularly important.
A duty attached to a permission through \texttt{odrl:duty} is a prerequisite
for that permission. A duty connected directly to a policy through
\texttt{odrl:obligation} is an independent obligation. Prior written consent
and a later annual review should therefore not be represented as an
undifferentiated list of actions that must all have happened before the first
transaction. The circular and the chosen event interpretation must determine
which performance is needed, by whom and at what time
\citep[secs.~2.6.1, 2.6.4--2.6.5]{odrlmodel2018}.

A \emph{profile} defines additional vocabulary and semantics for a particular
application. The proposed HKMA profile would name applying the Streamlined
Approach, reviewing continued qualification, providing a requested explanation,
and the observations used to establish their conditions. ODRL requires a
processor to understand a declared profile and stop when it does not recognise
it. A new identifier in a JSON file is consequently not enough: the evaluator
needs an implementation of that term's meaning. A profile extends the core;
it is not permission to reverse the meaning of an existing ODRL relationship
\citep[sec.~3]{odrlmodel2018}.

For this case, the resource can be the identified service or procedure instance
to which streamlined treatment applies. The intermediary is the party exercising
that permission; the client and account remain separately identified. Regulatory
source authority should be recorded explicitly in provenance. It should not
be inferred merely by putting the regulator into an \texttt{assigner} field,
which would hide a substantive mapping choice. A general regulatory rule set
can use \texttt{odrl:Set}; a particular agreement with a client is a different
record. These are proposed domain choices to specify and test, rather than
universal mappings supplied by ODRL.

\subsection{The circular contains calculations, dispensations and legal changes}

The financial limb illustrates the calculational part. Let $P$ be the
qualifying portfolio value and $N$ the qualifying net assets excluding the
primary residence, both expressed in Hong Kong dollars at the declared
assessment date. After the ownership, aggregation and currency questions have
been resolved, the financial condition is
\begin{equation}
\label{eq:odrl-financial-limb}
 F(P,N)=(P\geq 40{,}000{,}000)\ \lor\ (N\geq 80{,}000{,}000).
\end{equation}
Equality satisfies either threshold. A portfolio of HK\$40 million therefore
passes this limb even when the net-asset alternative does not. The disjunction
does not establish SPI status: professional-investor status, knowledge or
experience and investment objectives are additional conditions. Joint holdings,
written allocations, liabilities and qualifying investment-holding corporations
require their own source-based treatment
\citep[paras.~1--5]{sfcspiannex1}.

Knowledge or experience can be established through the specified degree or
postgraduate diploma, a professional qualification in finance, at least one
year of relevant professional work in the financial sector, or at least five
transactions within the past three years in the same product category.
Reliance on trading experience therefore needs category identity and a calendar
reference date. A representative's expertise is not automatically the client's
expertise, although transactions through a representative or joint account may
contribute to the client's trading history. Unknown objectives cannot establish
that the client is nonconservative
\citep[paras.~4.1, 5.1 and 7.3]{sfcspiannex1}
\citep[FAQ~2]{sfcspiannex2}.

The procedure relaxations introduce a different kind of statement. Paragraphs
10 and 11 say that specified procedures are ``not required'' under the
applicable conditions. This is a dispensation from a requirement; it is not
a prohibition on performing that procedure. Nor is it a blanket exemption
from suitability obligations. The explanation relaxation has exceptions for
requests and material queries, and solicited or recommended transactions differ
from unsolicited complex-product transactions. A translation must preserve the
scope of the relieved duty and the duties that survive
\citep[paras.~10--11]{sfcspiannex1}.

Client choices also change what applies. Prior written agreement and the
specified explanation precede use of the approach. Paragraph~13.1(d) provides
for withdrawal; paragraphs~13.2(c)(iii) and 13.2(d)(iii) concern changes to
categories and thresholds; paragraph~14 requires annual review and reminders.
An old consent event cannot authorize indefinitely after withdrawal. Equally,
``annually'' does not choose a 365-day arithmetic convention. The selected
calendar and effective-time interpretation must be visible
\citep[paras.~13--14]{sfcspiannex1}.

Finally, the financial evidence requirement must not be strengthened
accidentally. FAQ~1 permits reliance on disclosed information after reasonable
know-your-client efforts in the stated circumstances; inconsistencies need
clarification. Recordkeeping and universally demanding documentary proof are
different requirements. These qualitative assessments need evidence and a
declared interpretation. Inventing a numeric cutoff for ``reasonably satisfied''
would make a program easier to run by changing its subject
\citep[FAQ~1]{sfcspiannex2}.

\subsection{What direct text-to-ODRL research demonstrates}

\citet{revelodrl2023} offer an inspectable translation route: a sentence is
parsed into Abstract Meaning Representation (AMR), converted into RDF, and
transduced into ODRL. AMR makes linguistic relationships explicit before
the policy representation is selected. This is a useful baseline for separating
linguistic analysis from policy construction. It does not remove the need to
resolve definitions, references and the legal effect of exceptions.

Their experiment contains approximately 100 simple deontic sentences, about
20\% handcrafted and the remainder generated from a grammar, plus ten
negative examples. Modality extraction has reported precision 0.780, recall
0.640 and F-score 0.703; action extraction has precision 0.989, recall 0.810
and F-score 0.891. These are component scores rather than whole-policy legal
accuracy. Document preprocessing and definition resolution are described as
hypothetical extensions. The error discussion includes the scope of negation:
permission not to reproduce is different from lack of permission to reproduce.
That problem is directly relevant to the circular's ``not required'' language
\citep[method, experiments and error discussion]{revelodrl2023}.

\citet{mustafaodrl2024} take an ontology-guided language-model approach. They
provide ODRL concepts, curated explanations and examples, then refine generated
policies. Twelve cultural-dataspace use cases, three models and three methods
give 108 configurations; the paper describes initially running each use case
once. Its prominent 91.95\% figure is 217 of 236 selected checks for one use
case with refinement. It is not the probability that a complete legal policy
was interpreted correctly. The method supports a small generation baseline,
especially because it makes the allowed vocabulary explicit
\citep[secs.~2--4 and Table~1]{mustafaodrl2024}.

The associated LLM4ODRL code parses policies and checks SHACL shapes. SHACL
expresses conditions on an RDF graph: for example, whether a particular kind
of node has the required properties. Some repository shapes require creator,
description, title and issue date. Those can be useful application requirements,
but should not all be attributed to the W3C core standard. No repository-level
licence was found in the inspected checkout, so the method can inform a baseline
without assuming permission to copy its code \citep{hkmaodrlsoftware2026}.

\citet{agentodrl2026} add an orchestrator, reference rewriting, decomposition,
ontology-guided generation and SHACL-guided correction. Seventy seed cases are
expanded with a model into 770 cases: 440 simple, 220 with parallel relationships
and 110 with sequential dependencies. The semantic evaluation uses checkpoints
extracted by a model and judgments from two language models; the grammar score
uses selected shape violations. These measurements are useful for the authors'
generation experiment. They are not an independently established legal
interpretation of a circular and its annexes.

AgentODRL also gives a reason to keep the first experiment simple. Fixed
preprocessing workflows achieve higher aggregate semantic scores than the
adaptive orchestrator in parts of the reported comparison. The evidence does
not justify assuming that more orchestration improves this task. Compare a
direct ontology-guided proposal with explicit reference handling before adopting
the full system. Its Apache-licensed code is reusable, but the inspected
correction loop can retain a final candidate after exhausting its correction
budget without establishing success. The local result must preserve that failure
status. Reference rewriting must also retain the passages that supplied each
condition and exception \citep{agentodrl2026,hkmaodrlsoftware2026}.

There is earlier NLP work on RDF licence generation by Cabrio, Aprosio and
Villata, discussed in the AMR paper. Its original technical text was not
obtained in this survey, so it is a historical lead rather than separately
verified performance evidence. Likewise, De Vos, Kirrane, Padget and Satoh's
2019 \emph{ODRL Policy Modelling and Compliance Checking} is a directly
relevant lead on regulatory profiles and answer-set-program compliance,
including GDPR. The publication record and later references were located,
but the original attachment remained unavailable. Detailed claims about its
encoding or adoption would require that text
\citep{revelodrl2023,hkmaodrlsurvey2026}.

\subsection{Legal meaning and formal semantics constrain the translation}

\citet{criticalodrl2021} explain why a legal relationship can exceed a simple
usage permission. Duties, institutional powers, delegation, revocation and
policy lifecycle require distinctions that are easily hidden by a single
allow/deny result. For this circular, withdrawing consent changes the
arrangement; a dispensation changes which procedure is required; and the
regulator, intermediary and client occupy different roles. The critique
identifies modelling questions. Later semantics work must be considered
alongside it rather than assuming ODRL has remained unchanged.

The inspected AIUP profile supplies a practical example of regulatory
specialisation. It defines vocabulary for intended and precluded AI uses,
oversight and reporting, while expressly limiting its coverage of the EU AI
Act. Its draft specification, rather than the associated research paper, was
the source inspected here. Its relevance is the discipline of declaring exact
domain terms and supported questions. A corresponding HKMA profile should
state which duties and changes it can represent, rather than call every
operation \texttt{odrl:use} \citep{aiupprofile2024}.

\citet{bonattiodrl2025} develop declarative semantics for rule activation,
action matching, duties and compliance over traces. A trace records relevant
events and their order. This makes an otherwise hidden question explicit:
does one performance of a duty permit one subsequent action or many? A consent
to a continuing arrangement and an explanation in response to a particular
request require different answers. Their discussion also distinguishes
constraints on activation from conditions of fulfilment, considers policy gaps,
and identifies unresolved features such as temporal operands and a separate
meaning for \texttt{andSequence}. These are proposed semantic choices, not
a new W3C Recommendation
\citep[sec.~3, especially 3.6--3.8]{bonattiodrl2025}.

\citet{odrlnormalisation2026} show how a restricted policy fragment can be
normalised and compared under declared semantics. Their algorithms partition
numerical and symbolic conditions into matching regions and establish
preservation results. This suggests future checks for overlapping threshold
conditions or an unintended change between exports. The fragment excludes
attached duties, remedies and consequences; standalone obligations remain
unchanged in the policy-preservation result. Splitting one obligation into
several required obligations can change its meaning. The normal form can grow
exponentially in the number of attributes, and empirical performance evaluation
is left for future work. This is a basis for restricted comparison, not a
generic equivalence procedure for the whole circular
\citep[secs.~3--4, especially Theorem~6]{odrlnormalisation2026}.

LegalRuleML and eFLINT help identify what the internal interpretation should
retain even when the initial export is narrower. LegalRuleML addresses source
correspondence, context, defeasibility, authority and alternative readings.
The original eFLINT work separates action compliance from duty compliance
and represents events that create or terminate normative relationships.
Thus a transaction decision and an outstanding review duty can be reported
separately. These ideas support the internal design; this proposal does not
make a new LegalRuleML serializer or eFLINT runtime a prerequisite
\citep{legalruleml2021,eflint2020}.

\subsection{Evaluation software: what to reuse and what to test first}

\citet{oval2026} formalise access decisions and offline and streaming monitoring,
and present the OVAL implementation. Its algorithm processes events while
tracking rule matches and outstanding requirements. That is closer to the
circular's combination of procedure eligibility and continuing duties than
syntax validation alone. The reported evaluation uses synthetic workloads and
comparison with the SolidLab/FORCE cases. Agreement is 54 of 68 cases after
mapping the comparison system's ``not enough information'' to false. Different
reasoning and duty interpretations account for disagreements. Runtime results
on these workloads therefore cannot establish greater legal correctness
\citep[secs.~4--6]{oval2026}.

The pinned October OVAL code is a reasonable first candidate for a bounded
trial. Its access API defaults to accepting unfulfilled duties on a promise of
future fulfilment; another setting requires prior performance. The default
must not decide whether written consent can be obtained after the transaction.
Missing values can become false or remain null depending on configuration;
\texttt{andSequence} is implemented as conjunction rather than temporal
sequence; event counting is not a category-specific rolling three-year
history. Numeric comparisons attempt floating-point conversion, even where
older feature documentation describes lexical comparison. Explicit units,
some temporal operands and other constructs are unsupported. The proposal
therefore supplies exact calculations and a declared event adapter, with
unsupported terms reported before evaluation \citep{hkmaodrlsoftware2026}.

FORCE offers a useful testing pattern: a policy, request, state of the world
and expected compliance report. Its current documentation describes limitations
on duties, while code includes duty-reporting structures and rules. Reporting
a duty and enforcing it as a prerequisite to permission are distinct
capabilities. A helper also connects policy-level duties through
\texttt{odrl:duty}, where the W3C model uses \texttt{odrl:obligation}.
These are concrete compatibility checks for the intended input route, not
evidence that every FORCE evaluation fails. The reported 67/68
activation-similarity result is not a local result or a measure of fulfilled
legal duties. The associated 2025 papers were located bibliographically but
not obtained in full, so the recommendation rests on inspected code and
documentation \citep{hkmaodrlsoftware2026}.

ODRE addresses the connection between descriptive policy terms and executable
values. \citet{odre2024} use runtime interpolation and operand/action adapters,
with Python and Java implementations and 24 demonstration tests per
implementation. Their tested enforcement is synchronous; temporal monitoring
is future work in that version. The Python checkout presents additional static
concerns: the rule traversal discards whether a rule was reached through
permission, prohibition or obligation, and the constraint query appears to
look directly on a rule where its own sample uses a nested constraint node.
An empty constraint result is then accepted. These findings were not executed
as failures, but they are sufficient to withhold selection as the first general
normative evaluator pending focused reproduction
\citep{odre2024,hkmaodrlsoftware2026}.

Eclipse Dataspace Components provides registered evaluation functions,
execution scopes and JSON-LD policy transformation. These are useful later
if the product needs dataspace negotiation or transfer services. They do not
extract a circular's norms or define SPI qualification. pySHACL is the smaller
immediate dependency for graph validation; its shapes must be reviewed and
must actually target the intended policy nodes. The public ODRL-shape
repository is a source of ideas, but no repository-level licence was found
in the inspected checkout \citep{hkmaodrlsoftware2026}.

\begin{table}[htbp]
\centering
\caption{Selection for a first trial, based on static inspection rather than a benchmark ranking.}
\label{tab:odrl-software-choices}
\begingroup
\footnotesize
\renewcommand{\arraystretch}{1.08}
\begin{tabularx}{\textwidth}{P{27mm}YY}
\toprule
Component & Proposed use & Condition before reliance\\
\midrule
\href{https://github.com/keman-ruc/wanle-Agentodrl}{AgentODRL} & Generation and reference-handling ideas & Preserve source maps and failed correction status.\\
\href{https://github.com/Daham-Mustaf/LLM4ODRL}{LLM4ODRL} & Simple ontology-guided methodological baseline & Resolve licence before copying code; review shape assumptions.\\
\href{https://github.com/RDFLib/pySHACL}{pySHACL} & Structural and profile checks & Review shapes, targeting and negative cases.\\
\href{https://github.com/DIPS-Tools/ODRL-Engine}{OVAL} & First evaluator trial & Pin duty/default semantics; supply exact facts and explicit time handling.\\
\href{https://github.com/odrl-force/ODRL-Evaluator}{FORCE} & Scenario/report design and second evaluator & Establish duty and policy-obligation behavior on the chosen route.\\
\href{https://github.com/ODRE-Framework/odre-python}{ODRE} & Operand-adapter design reference & Reproduce and resolve the inspected traversal concerns.\\
\href{https://github.com/eclipse-edc/Connector}{EDC} & Later deployment integration & Supply the domain semantics; justify the larger integration cost.\\
\bottomrule
\end{tabularx}
\endgroup
\end{table}

The inspected AgentODRL, pySHACL documentation, ODRE and EDC declare Apache-2.0
licensing; OVAL and FORCE use MIT. The table links to the public repositories.
The proposed first OVAL trial is pinned to \texttt{b2838782b7ce}; the maintained
AgentODRL checkout is \texttt{d55ed231ba7d}, FORCE is \texttt{194894a04588},
and ODRE Python is \texttt{1a2cb06004f0}. All nine full checkout commits and inspected-file
hashes are retained in the survey manifest. These snapshots fix what was
examined, not a promise about a later release. No software is ranked as
production ready on this evidence \citep{hkmaodrlsoftware2026}.

\subsection{Choosing the relationship with Catala}

Catala is valuable where the circular gives computations with defaults and
exceptions. Equation~\ref{eq:odrl-financial-limb}, explicit allocation rules
and category histories are candidates once their premises are fixed. The
Catala paper's compilation result concerns translation between formal calculi.
It does not prove that a natural-language source was interpreted correctly
or that another language has the same semantics
\citep[secs.~4.2--4.5]{catala2021}.

Translating the document into Catala first gives early executable calculations
and uses existing tooling. It is attractive for a predominantly computational
source or an already reviewed Catala model. Here, however, a calculation need
not retain the distinction between a dispensation, a legal power to withdraw,
and an outstanding duty. Reconstructing that information later for ODRL can
lose meaning. Catala-first is therefore an unnecessary prerequisite for this
circular.

ODRL-first followed by Catala delivers the desired interchange representation
early and makes parties, actions and duties visible. But any meaning omitted
from ODRL is also missing from the subsequent calculation. Agreement between
the two generated outputs mainly tests that translation step. Fully separate
document-to-ODRL and document-to-Catala processes provide more independence,
but duplicate interpretation work and permit sources, terminology and
assumptions to drift. They are useful experimental comparators on a small
fragment, rather than the preferred everyday authoring process.

The recommended design retains a richer reviewed interpretation and exports
ODRL while invoking Catala for selected calculations. Independent scenario
specifications challenge both. This separates four questions: whether the
graph conforms to its profile, whether the evaluator implements the chosen
semantics, whether the two computations agree on a declared domain, and
whether the interpretation is supported by the circular. Catala can help
answer the third and expose defects relevant to the second. It does not
establish the first or fourth by itself.

The practical consequence is to begin with source coverage and the FAQ~6,
consent and explanation counterexamples. The initial feature should produce a
reviewable policy and an explicit account of what it cannot express. A generic
ODRL-to-Catala compiler is not required to obtain that result.
Section~\ref{sec:odrl-proposal} specifies the records, modules and checks
needed to make this recommendation concrete.
